\cotefile{45}
\scanpage{43}
\mspage{23}
\begin{center}\uuline{BIBLIOGRAPHIE}\end{center}

[1] C.ARF. \uline{Untersuchungen über reinverzweigte Erweiterungen diskret bewerteter perfekter Körper}. Journ. Crelle, 181, 1940, p.1-44.

[2] E. ARTIN. \uline{Die gruppentheoretische Struktur der Diskriminante algebraischer Zahlkörper}. Journ. Crelle, 164, 1931, p.1-11.

[3] R. BRAUER. \uline{On the algebraic structure of group rings}. Journ. Soc. Math. Japan, 3, 1951, p.237-251.

[4] C. CHEVALLEY et A.WEIL. \uline{Über das Verhalten der Integrale} 1 \uline{Gattung bei Automorphismen des Funktionenkörpers}. Hamburg Abh., 10, 1934, p.358-361.

[5] M. DEURING. \uline{Algebren}. Ergebnisse der Math., n°IV-1, Springer, 1935.

[6] M. DEURING. \uline{Die Typen der Multiplikatorenringe elliptischer Funktionenkörper}. Hamburg Abh., 14, 1941, p.197-272.

[7] H. HASSE. \uline{Normenresttheorie galoisscher Zahlkörper mit Anwendungen auf Führer und Diskriminante abelscher Zahldörper}. Journ. Fac. Sci. Tokyo, 2, 1934, p.477-498.

[8] J.HERBRAND. \uline{Sur la théorie des groupes de décomposition, d'inertie, et de ramification}. Journ. Liouville, 10, 1931, p.481-498.

[9] S.LANG. \uline{Abelian varieties}. Interscience Tracts n°7, New-York, 1959.

[10] P.ROQUETTE. \uline{Realisierung von Darstellungen endlicher nilpotenter Gruppen}. Archiv der Math., 9, 1958, p.241-250.

[11] P.SAMUEL et O.ZARISKI. \uline{Commutative Algebra} I. Van Nostrand, 1958.

[12] O.SCHILLING. \uline{Über die Darstellungen endlicher Gruppen}. Journ. Crelle, 174, 1936, p.18.

[13] J.-P.SERRE. \uline{Corps locaux et isogénies}. Sém.BOURBAKI, 1959, n°185.

[14] R.SWAN. \uline{Induced representations and projective modules}. A paraître prochainement.

[15] A.WEIL. \uline{L'avenir des mathématiques}. Cahiers du Sud, Marseille, 1947, ainsi que Bol. S\={a}o Paulo, 1, 1946, p.55-68.

[16] A.WEIL. \uline{Sur les courbes algébriques et les variétés qui s'en déduisent}. Paris, Hermann, 1948.

[17] E.WITT. \uline{Schiefkörper über diskret bewerteten Körpern}. J. Crelle, 176, 1936, p.153-156.

[18] E.WITT. \uline{Die algebraische Struktur des Gruppenringes einer endlichen Gruppe über einem Zahlkörper}. Journ.Crelle, 190, 1952, p.231-245.

\begin{center}***\end{center}
% archive sheet number: 42
